\chapter{Conclusion}\label{ch_conclusion}

The experiments in the preceding chapter evaluated each optimisation level in isolation and then measured their combined effect, revealing that the cascade from style simplification through tile shaving to columnar re-encoding is more effective than applying either level in isolation.
This chapter draws together the findings: it summarises the contributions, distils the key insights, revisits the three research questions posed in \cref{sec_research_goals}, and identifies directions for future work.

\section{Summary of Contributions}

The main contributions are:

\begin{enumerate}
    \item \textbf{Two-phase style optimisations} (47 expression-level rewrites, 6 structural passes).
    Expression simplification proved most valuable as an \emph{enabler}: constant folding and algebraic simplification reduced filters to forms that dead elimination and layer merging could exploit, producing the largest rendering gains by cutting per-frame layer and draw-call counts.

    \item \textbf{Style-driven tile shaving} demonstrated a cascade effect: each style layer removed by dead elimination shrank the pruning advisory, which in turn eliminated entire columns and tiles from the encoded output.

    \item \textbf{Automatic encoding strategy selection} for the \ac{MLT} format reduced tile volume by \qty{28}{\percent} over the fixed-plan \ac{MLT}-Java baseline and by \qty{12}{\percent} over gzip-compressed \ac{MVT}, with the advantage preserved even against zstd-22.

    \item \textbf{Empirical evaluation} across 15 styles, 18 scenarios, and 17 cumulative ablation steps.
    The full pipeline reduced total tile volume by \qty{15}{\percent} (\qty{3.07}{\giga\byte} $\to$ \qty{2.61}{\giga\byte}), improved load times by \qtyrange{1}{33}{\percent} for most scenarios, and held FPS stability within \qty{\pm 1}{\percent}.
\end{enumerate}

\section{Key Findings}

Three architectural insights emerged from the evaluation that go beyond the quantitative results reported in \cref{ch_experiments}.

\emph{First}, the cascade effect --- style simplification makes data optimisations more effective --- is the central finding.
Each style layer removed by dead elimination shrinks the pruning advisory, which eliminates entire property columns and even whole source-layers from the encoded tile.
The pipeline is more effective than applying either level in isolation, and any evaluation that benchmarks data-level or style-level optimisations alone will understate the combined benefit.
For the geospatial community this means that style authoring and tile encoding should no longer be treated as independent concerns: the two interact, and tooling that exploits the interaction can recover substantial transfer and rendering savings without any change to the underlying data or the renderer.

\emph{Second}, different optimisation categories address different bottlenecks in the rendering pipeline.
Structural passes (dead elimination, layer merging) reduce per-frame draw-call overhead and dominate the rendering-throughput improvements.
Data-level passes (tile shaving, columnar re-encoding) reduce transfer size and dominate the load-time improvements.
Expression passes serve primarily as \emph{enablers}, simplifying filters into forms that structural and data-level passes can exploit.
This decomposition matters for deployment decisions: a mobile deployment on a constrained network should prioritise tile shaving and re-encoding, whereas a desktop deployment bottlenecked by draw calls benefits most from structural passes alone.

\emph{Third}, the compiler-optimisation framing --- peephole rewriting, fixpoint iteration, dead code elimination --- transfers effectively to the map-style domain.
The 47 rewrite rules are individually simple, yet their composition through fixpoint iteration achieves globally significant simplification across styles of widely varying complexity (raw reductions of \qtyrange{10}{71}{\percent} across the Maputnik generalisation corpus).
This suggests that other declarative rendering specifications (CSS, SVG filters, shader graphs) may benefit from the same approach: a small rule set applied to fixpoint, followed by structural passes that exploit the simplified form.

\section{Research Questions Revisited}\label{sec_rq_revisited}

The three research questions posed in \cref{sec_research_goals} can now be answered directly from the experiments in \cref{ch_experiments}.

\paragraph{R1 - To what extent can automatic data and style co-optimisations reduce map tile size and improve rendering performance?}
Style-level evaluation across all 15 styles and 18 scenarios shows style-gzip reductions of \qtyrange{3}{47}{\percent} (median ${\sim}$\qty{8}{\percent}) and \ac{FPS} stability within \qty{\pm 1}{\percent} of baseline.
At the tile-data level, full-pipeline benchmarks on the two styles evaluated end-to-end (Fiord and Liberty) show a total tile volume reduction of roughly \qty{15}{\percent} (from \qty{3.07}{\giga\byte} to \qty{2.61}{\giga\byte} on the Germany OpenMapTiles tileset) and load-time improvements of \qtyrange{71}{75}{\percent} (\cref{sub_combined_ablation}).
The gains are not uniform: verbose institutional styles with hundreds of layers benefit far more than already-compact basemaps, and scenarios dominated by overview zoom levels benefit more than those dominated by street-level panning (\cref{sub_combined_marginal}).
The quantitative answer to R1 is therefore \enquote{yes, meaningfully, but in a style- and scenario-dependent way}, with the style-driven cascade into tile shaving and the Rust encoder's strategy competition accounting for most of the tile-size reduction, and dead elimination plus layer merging accounting for most of the frame-rate and load-time reductions.

\paragraph{R2 - Which techniques are most effective, and how do they interact?}
The cumulative ablation in \cref{sub_combined_ablation} and the per-style results in \cref{app_per_style} rank the passes by marginal contribution and show three clusters.
\emph{First}, constant folding (including stats-driven folding) and expression simplification dominate the expression-level reductions because they simplify filters into forms that the structural passes can exploit, and because a single successful fold cascades through enclosing expressions.
\emph{Second}, dead elimination and layer merging dominate the structural-level reductions because they reduce the per-frame draw-call count directly; they are also strongly synergistic, because dead elimination removes gaps between mergeable layers that would otherwise prevent the merge pass from recognising adjacency.
\emph{Third}, data-level optimisations operate on a different axis: tile shaving reduces transfer size by cascading the optimised style's property and value references into a pruning advisory, and the \ac{MLT} encoder's sort-strategy and per-stream competition reduces encoded size on top of that.
The three clusters interact through the \emph{cascade effect} documented in \cref{sub_mlt_results}: style-level dead elimination shrinks the pruning advisory, which shrinks the tiles, which shrinks both the transfer time and the decoder's work - so evaluating any one category in isolation understates its contribution.

\paragraph{R3 - How expensive are these optimisations in terms of preprocessing time, and what are the tradeoffs between transfer size, decoding latency, and rendering throughput?}
The preprocessing-cost model in \cref{sub_preprocessing_cost} decomposes the total cost into four additive stages (style optimisations, statistics collection, advisory computation, \ac{MLT} re-encoding) and shows that the dominant term is re-encoding.
The multi-strategy competition increases re-encoding cost by roughly a factor of three to five over a fixed-plan encoder, but this cost is paid once per tileset build and amortised over every subsequent tile request; the \ac{RAII} alternatives guard, warm-start \ac{FSST} caching, and scratch-buffer recycling keep the absolute cost of each additional trial low enough that the full pipeline still runs in tens of minutes on the Germany OpenMapTiles tileset on a single-threaded desktop workstation.
On the tradeoff axis, transfer-size reduction is the most robust gain (style-level gzip reductions of \qtyrange{3}{47}{\percent}, median ${\sim}$\qty{8}{\percent}; tile-level volume reduction of \qty{12.4}{\percent} for \ac{MLT}-Rust vs.\ \ac{MVT} under gzip).
Decoding latency, measured via Criterion micro-benchmarks on the OpenMapTiles fixtures (\cref{tab:decode_throughput}), shows that the \ac{MLT}-Rust decoder achieves \numrange{4.6}{8.3}$\times$ higher full-decode throughput than \ac{MVT}+gzip across zoom levels 4, 7, and 13.
Rendering throughput (FPS) is reported in \cref{sub_combined_ablation} and shows \qty{\pm 1}{\percent} stability, affected almost entirely by the structural passes rather than by the data-level passes --- confirming that data-style co-optimisation is not a single knob but a set of complementary mechanisms that each address a different bottleneck.
Energy consumption is not measured directly (the benchmark host is a desktop workstation rather than a battery-constrained device), but the transfer-size and frame-rate results are well-established proxies on mobile targets and the qualitative discussion in \cref{sub_combined_discussion} translates them into the expected battery-life direction.

\section{Future Work}\label{sec_conclusion_future_work}

Several directions remain open:

\begin{itemize}
    \item \textbf{Symbol layer merging}: the current implementation excludes symbol layers from merging because it would disrupt collision detection. Investigating collision-aware merging strategies could yield additional layer count reductions on label-heavy styles.
    \item \textbf{Runtime re-encoding}: the current tile encoding pipeline is offline. Exploring on-the-fly MVT-to-MLT re-encoding at the tile server or \ac{CDN} edge could make the encoding optimisations available without a full tile rebuild.
    \item \textbf{Multi-objective optimisations}: the post-hoc Pareto analysis in \cref{sec_exp_pareto} reports a two-axis size-cost frontier over the measured pass configurations (cumulative trajectory plus isolated mode).
    Its deliberate scope is the result of a noise-floor decomposition that showed the timing-derived axes (\texttt{loadMs}, \texttt{meanFrameMs}, \texttt{peakHeapMB}, and a CPU-time energy proxy) could not support per-step frontier claims at the current sample size.
    Extending the frontier to those axes requires either a substantial increase in run count per cell (on the order of $5$-$10\times$), a dedicated benchmark host with reduced thermal variability, or renderer-internal instrumentation that isolates decoder time from the surrounding fetch-and-paint work - and most likely some combination of all three.
    Beyond reducing noise, a formal \ac{MOOP} search over the full pass-permutation and internal-knob space would extend the analysis layer with a true search procedure rather than a reading of already-measured data, and is out of scope for this thesis.
    \item \textbf{Broader style coverage}: extending the evaluation to proprietary and commercial styles (Mapbox Streets, Google Maps style equivalents) would strengthen the generalisability claims.
    \item \textbf{Client-side integration}: integrating the optimised \ac{MLT} decoder into MapLibre GL~JS and measuring end-to-end improvements in a production deployment.
    \item \textbf{Serving-time applicability}: the current pipeline assumes offline batch processing, which is impractical for dynamic tile sources (e.g., PostGIS-backed tile servers with frequently updated data).
    Not all passes require offline processing, however.
    The style optimiser operates on the style alone and can run at deployment time (when the style is published) rather than at tile-build time.
    The tile pruning advisory could potentially be applied as a lightweight per-request filter at serving time, since it reduces to column selection and predicate evaluation --- operations that tile servers already perform for other purposes.
    The encoding strategy competition, by contrast, requires access to the full column to profile and is inherently a batch operation; its results could be cached per schema and reused across tile requests.
    Characterising which passes are applicable at serving time and measuring their overhead on a request-critical path would extend the pipeline's applicability to dynamic tile sources.
\end{itemize}
